\documentclass[10pt,a4paper]{amsart}
\usepackage[margin=19mm]{geometry}
\usepackage{amsmath,amssymb,mathtools}
\usepackage[scr=rsfs,cal=euler]{mathalfa}
\usepackage{xfrac}
\usepackage[unicode,hidelinks,bookmarks=false]{hyperref}
\newcommand{\ie}{i.e.\ }
\newcommand{\cf}{cf.\ }
\newcommand{\resp}{respectively\ }
\newcommand{\Subsection}{\subsection}
\providecommand{\nobreakdash}{\nobreak\hbox{-}}
\setlength{\emergencystretch}{2em}
\multlinegap0pt
\usepackage{bm}
\usepackage{bbm}
\usepackage{dsfont}
\usepackage{xspace}
\usepackage{cancel}
\usepackage{mathtools}
\usepackage{amsthm}
\makeatletter
\@ifundefined{lem}{\newtheorem{lem}{Lemma}}{}
\makeatother
\newcommand{\dd}{\mathop{}\!\mathrm{d}}
\title[Estimates for the Gross-Pitaevskii equation linearized aroun]{Estimates for the Gross-Pitaevskii equation linearized around a vortex}
\author{Charles Collot, Pierre Germain, Eliot Pacherie}
\date{Source: arXiv:2503.02953v3 (CC BY 4.0)}
\begin{document}
\maketitle

\noindent\emph{Source and licence.} Estimates for the Gross-Pitaevskii equation linearized around a vortex. Version arXiv:2503.02953v3 (CC BY 4.0), distributed under the Creative Commons Attribution 4.0 International licence, as recorded in the arXiv licence metadata for this exact version and shown on the abstract page https://arxiv.org/abs/2503.02953v3. Full rights notice: licenses/arxiv-2503.02953-CC-BY-4.0.txt.\par\smallskip

\noindent\textit{Editorial scope.} Counted unit: the Lem whose source label is \texttt{latecheckout} in the article (source0.tex lines 5933--5968, source label \texttt{latecheckout}), delivered together with the article's own complete original proof of it (source0.tex lines 5970--6321, one \texttt{proof} environment of 352 lines). No mathematics is written, repaired or re-derived: statement and proof are the article's own text line for line. Required context retained uncounted: maineq3 (lines 3134--3142); lydmor (lines 3435--3444); factexp (lines 3545--3588); lethargy (lines 4048--4059); kast (lines 5202--5218); eqsec9 (lines 5240--5318); alegant (lines 5504--5536); heimat (lines 5893--5897). Every definition, display and label of those lines is retained unchanged. Omitted from this package and not redistributed: 1--3133, 3143--3434, 3445--3544, 3589--4047, 4060--5201, 5219--5239, 5319--5503, 5537--5892, 5898--5932, 5969--5969, 6322--8585. Nothing inside a retained excerpt is omitted or altered: every retained body is delivered exactly as the article's lines are written, so no omission receipt of any kind accompanies this package. Citations: the retained excerpts carry no citation command at all, so no bibliography entry of the article is transcribed and no reference is redistributed. Reference adaptation: none was needed and none was made; every reference inside the delivered unit resolves inside it, so no unresolved reference or citation is produced. Printed slips kept literal: no printed word, symbol, hypothesis or number of the counted statement or of its proof was altered. Layout and class: the article's own class is \texttt{amsart} with packages including inputenc, tikz, tikz-3dplot, amsmath, amsthm, amssymb, graphics, esint, amscd, color, xcolor, mathtools, url, bbm; the wrapper is the lane's pinned \texttt{amsart} editorial wrapper at 10pt on A4, which restates the theorem-like environments the retained text needs and adds the article's own macro definitions that the retained text needs (\texttt{\textbackslash dd}), with extra wrapper packages (bm, bbm, dsfont, xspace, cancel, mathtools). The delivered numbering is the unit's own. Assets: no retained span contains a figure environment, a graphics-inclusion command, a PDF-page inclusion, a table or a TikZ picture; no image file is copied into this package, the package declares no asset dependency and no image file is redistributed with it. Licence: version arXiv:2503.02953v3 of this article is distributed under the Creative Commons Attribution 4.0 International licence, as recorded in the arXiv licence metadata for this exact version and shown on the abstract page https://arxiv.org/abs/2503.02953v3. A full rights notice is delivered at licenses/arxiv-2503.02953-CC-BY-4.0.txt. Characters: every retained excerpt is pure ASCII, so the delivered engine, XeLaTeX with its default Latin Modern text font, reports no missing character.
  \begin{equation}
    \left(\begin{array}{c}
      (\Delta - \kappa^2) \varphi\\
      (\Delta + \xi^2) \psi
    \end{array}\right) + V_{\lambda} \left(\begin{array}{c}
      \varphi\\
      \psi
    \end{array}\right) = 0 \label{maineq3}
  \end{equation}

\begin{equation}
  \left(\begin{array}{c}
    \left( \Delta - \left( \frac{\kappa}{\xi} \right)^2 -
    \frac{1}{y^2} \right) \varphi\\
    \left( \Delta + 1 - \frac{1}{y^2} \right) \psi
  \end{array}\right) (y) = \widetilde{V}_{\lambda} \left(\begin{array}{c}
    \varphi\\
    \psi
  \end{array}\right) (y) \label{lydmor}
\end{equation}

\begin{lem}
  \label{factexp}For any $\lambda > 0$, there exists a solution of
  (\ref{lydmor}) of the form
  \[ \left(\begin{array}{c}
       \varphi_2\\
       \psi_2
     \end{array}\right) =  K_1 \left( \frac{\kappa}{\xi} \cdot \right)
     \left(\begin{array}{c}
       h_1\\
       h_2
     \end{array}\right) \]
such that, for any $\lambda_0 > 0$, taking $\lambda \geqslant \lambda_0$ and
  $y \geqslant 1$, we have for any $j, k \in \mathbb{N}$,
  \[ \left| \partial_{\lambda}^j \partial_y^k \left( \left(\begin{array}{c}
       h_1\\
       h_2
     \end{array}\right) - \left(\begin{array}{c}
       1\\
       0
     \end{array}\right) \right) \right| \lesssim_{\lambda_0, j, k}
     \frac{1}{(\xi + y) \lambda^j y^k} . \]
     
  Furthermore, given $y_0 \geqslant 1$, we have
  \[ \left(\begin{array}{c}
       \varphi_2\\
       \varphi_2'\\
       \psi_2\\
       \psi_2'
     \end{array}\right) (y_0) =   
     \left(\begin{array}{c}
       K_1 \left(y_0 \right)\\
       K_1' \left(y_0 \right)\\
       0\\
       0
     \end{array}\right) + o_{\lambda \rightarrow + \infty} (1)  \]
and for any $k \in \mathbb{N}^{\ast}$,
  \[ \left| \partial_{\lambda}^k \left(\begin{array}{c}
       \varphi_2\\
       \varphi_2'\\
       \psi_2\\
       \psi_2'
     \end{array}\right) (y_0) \right| \lesssim_{\lambda_0, y_0, k}
     \frac{1}{\xi \lambda^k} . \]
\end{lem}

\begin{lem}
  \label{lethargy}There exist two linearly independent solutions of $(L_0 - 2
  \rho^2) (Q) = 0$ denoted $Q_0$ and $\widetilde{Q}_0$, and they satisfy
  \[ Q_0 (r) \approx r,\qquad \widetilde{Q}_0 (r) \approx \frac{1}{2 r} \]
  when $r \rightarrow 0$, and
  \[ Q_0 (r) \approx \frac{C_1}{\sqrt{r}} e^{\sqrt{2} r},\qquad \widetilde{Q}_0 (r)
     \approx \frac{C_2}{\sqrt{r}} e^{- \sqrt{2} r} \]
  \[ \  \]
  when $r \rightarrow + \infty$ for some constants $C_1, C_2 > 0$.
  Furthermore, both $Q_0$ and $\widetilde{Q}_0$ are strictly positive functions on
  $\mathbb{R}^{+ \ast}$.
\end{lem}

\begin{lem}
  \label{kast}The quantities
  \[ a_{\pm} (\lambda) = \lambda \pm \langle \lambda \rangle, \; \xi =
     \sqrt{\langle \lambda \rangle - 1}, \; \kappa = \sqrt{1 + \langle \lambda
     \rangle} \]
  satisfy the following estimates for all $\lambda \in [0, 1]$ and $j \in \mathbb{N}$:
$$
| \partial_{\lambda}^j \xi | \lesssim_j \lambda^{\max (0, 1 - j)},
\qquad  | \partial_{\lambda}^j \langle \lambda \rangle | + | \partial_{\lambda}^j a_{\pm} | + | \partial_{\lambda}^j \kappa | \lesssim_j 1
$$
as well as
  \[ \left| \partial_{\lambda}^j \left( \frac{\langle \lambda \rangle -
     1}{\lambda^2} \right) \right| + \left| \partial_{\lambda}^j \left(
     \frac{\kappa^2 - 2}{\lambda^2} \right) \right| + \left|
     \partial_{\lambda}^j \left( \frac{\xi}{\lambda} \right) \right|
     \lesssim_j 1. \]
\end{lem}

\begin{lem}
  \label{eqsec9}
\begin{itemize}
\item[1)] (Perturbation of $\lambda = 0$ close to $r = + \infty$) Equation
  (\ref{maineq3}) is equivalent to
  \[ \left(\begin{array}{c}
       \left( \Delta - \kappa^2 + \frac{\kappa^2}{r^2} \right) \varphi\\
       (\Delta + \xi^2) \psi
     \end{array}\right) =\mathcal{V}_{\lambda} \left(\begin{array}{c}
       \varphi\\
       \psi
     \end{array}\right) \]
  where for all $j, k \in \mathbb{N}, r \geqslant y_0 \xi^{- 1}/2$
  \[ | \partial_{\lambda}^j \partial_r^k \mathcal{V}_{\lambda} | \lesssim_{j,
     k,y_0} \frac{\lambda^{\max (0, 1 - j)}}{(1 + r)^{2+k}} . \]
  
\item[2)] (At the scale of oscillations) In the variable $y = \xi r$, this equation
  can be written as
  \[ \left(\begin{array}{c}
       \left( \Delta_y - \left( \frac{\kappa}{\xi} \right)^2 + \frac{2}{y^2}
       \right) (\varphi)\\
       (\Delta_y + 1) (\psi)
     \end{array}\right) = \widetilde{M} \left(\begin{array}{c}
       \varphi\\
       \psi
     \end{array}\right) (y) \]
  where we have the estimate for any $j, k \in \mathbb{N}, \lambda \in [0, 1],
  y \geqslant y_0/2$,
  \[ | \partial_{\lambda}^j \partial_y^k \widetilde{M} (y, \lambda) | \lesssim_{j,
     k,y_0} \frac{\lambda^{- j}}{y^{2 + k}} \left(\begin{array}{cc}
       \lambda^2 & \lambda\\
       \lambda & \lambda^2
     \end{array}\right) \]
(where the bound is understood coordinatewise).
  
\item[3)] (Perturbation of the case $\lambda = 0$ for $r \lesssim \xi^{-1}$) The equation can also be written 
  \[ \left(\begin{array}{c}
       (L_0 - 2 \rho^2) (\varphi)\\
       L_0 (\psi)
     \end{array}\right) = \bar{V}_{\lambda} \left(\begin{array}{c}
       \varphi\\
       \psi
     \end{array}\right) \]
  where for any $j, k \in \mathbb{N}, r \geqslant 0$,
  \[ | \partial_{\lambda}^j \partial_r^k \bar{V}_{\lambda} | \lesssim_{j, k}
     \frac{\lambda^{\max (0, 2 - j)}}{(1 + r)^k} + \frac{\lambda^{\max (0, 1 -
     j)}}{(1 + r)^{2 + k}} \]
  and
  \[ \left| \partial_{\lambda}^j \partial_r^k \left(
     \frac{\bar{V}_{\lambda}}{\lambda} \right) \right| \lesssim_{j, k}
     \frac{1}{(1 + r)^k} . \]
  
\item[4)] (Combination of 1. and 3. where $L_0$ and $\xi^2$ are considered comparable) We
  also have the decomposition
  \[ \left(\begin{array}{c}
       (L_0 - 2 \rho^2) (\varphi)\\
       (L_0 + \xi^2) \psi
     \end{array}\right) = \widetilde{V}_{\lambda} \left(\begin{array}{c}
       \varphi\\
       \psi
     \end{array}\right) \]
  where
  \[ \widetilde{V}_{\lambda} = \left(\begin{array}{cc}
       \kappa^2 - 2 & - \lambda \frac{\rho^2 (r) - 1}{\langle \lambda
       \rangle}\\
       \frac{- \lambda}{\langle \lambda \rangle} (\rho^2 (r) - 1) & - 2 (1 -
       \langle \lambda \rangle) \frac{\rho^2 (r) - 1}{\langle \lambda \rangle}
     \end{array}\right) \]
  which satisfies for any $j, k \in \mathbb{N}$ that,
  \[ | \partial_{\lambda}^j \partial_r^k \widetilde{V}_{\lambda} | \lesssim_{j, k}
     \frac{1}{(1 + r)^{2 + k}} \left(\begin{array}{cc}
       0 & \lambda^{\max (0, 1 - j)}\\
       \lambda^{\max (0, 1 - j)} & \lambda^{\max (0, 2 - j)}
     \end{array}\right) + \delta_{k = 0} \left(\begin{array}{cc}
       \lambda^{\max (0, 2 - j)} & 0\\
       0 & 0
     \end{array}\right) . \]
\end{itemize}
\end{lem}

\begin{lem}
  \label{alegant} Given $\lambda$ small enough, the Jost solution decaying exponentially fast (that is the one constructed in Lemma \ref{factexp} and thus solution of (\ref{maineq3})) has the form
  \[ \left(\begin{array}{c}
       \varphi_2\\
       \psi_2
     \end{array}\right) (r) = K_i (\kappa r) \left(\begin{array}{c}
       h_1\\
       h_2
     \end{array}\right) (r) \]
  where for any $r \geqslant r_\ast, j, k \in \mathbb{N}$,
  \[ \left| \partial_{\lambda}^j \partial_r^k \left( \left(\begin{array}{c}
       h_1\\
       h_2
     \end{array}\right) - \left(\begin{array}{c}
       1\\
       0
     \end{array}\right) \right) \right| \lesssim_{j, k} \frac{\lambda^{\max
     (0, 1 - j)}}{(1 + r)^{1 + k}} . \]
  Furthermore, for any $j \in \mathbb{N}, \nu \in \{ \frac 12, 4 \},$
  \[ \left| \partial_{\lambda}^j \left( \frac{1}{K_i \left( y_0\nu
     \frac{\kappa}{\xi} \right)} \left(\begin{array}{c}
       \varphi_2\\
       \varphi'_2\\
       \psi_2\\
       \psi_2'
     \end{array}\right) (\nu y_0 \xi^{- 1}) - \left(\begin{array}{c}
       1\\
       - \sqrt{2}\\
       0\\
       0
     \end{array}\right) \right) \right| \lesssim_{j,y_0} \lambda^{\max (0, 2 - j)} .
  \]
\end{lem}

\begin{lem}
  \label{heimat}There exists $c_1 > 0$ such that
  \[ \frac{Q_0 (r)}{I_i \left( \sqrt{2} r \right)} \in c_1 +\mathcal{W}_2
     (r_\ast^+) . \]
\end{lem}

\begin{lem}
  \label{latecheckout}There exists $\lambda_0 > 0$ such that, for any $\lambda
  \in [0, \lambda_0]$, there exists a solution of (\ref{maineq3}) of the form
  \[ \left(\begin{array}{c}
       \Phi_1\\
       \Psi_1
     \end{array}\right) (r) = Q_0 (r) \left(\begin{array}{c}
       h_1\\
       h_2
     \end{array}\right) (r) \]
  with, for any $j, k \in \mathbb{N}, r \in [0, 8 y_0 \xi^{- 1}]$,
  \[ \left| \partial_{\lambda}^j \partial_r^k \left(
     \frac{1}{\lambda} \left[\left(\begin{array}{c}
       h_1\\
       h_2
     \end{array}\right) - \left(\begin{array}{c}
       1\\
       0
     \end{array}\right) \right] \right) \right| \lesssim_{j, k} (1 + r)^{j -
     k} . \]
  Furthermore, for any $j \in \mathbb{N} , \nu \in \{ \frac 12,4 \}$,
  \[ \left| \partial_{\lambda}^j \left( \frac{1}{I_i \left( \nu y_0
     \frac{\kappa}{\xi} \right)} \left(\begin{array}{c}
       \Phi_1\\
       \Phi_1'\\
       \Psi_1\\
       \Psi_1'
     \end{array}\right) (\nu y_0\xi^{- 1}) - \left(\begin{array}{c}
       1\\
       \sqrt{2}\\
       0\\
       0
     \end{array}\right) \right) \right| \lesssim_{j,y_0} \lambda^{2-j} .
  \]
  \[ \  \]
\end{lem}

\begin{proof}
  By Lemma \ref{eqsec9}.3, we use the equation
  \[ \left(\begin{array}{c}
       (L_0 - 2 \rho^2) (\varphi)\\
       L_0 (\psi)
     \end{array}\right) = \bar{V}_{\lambda} \left(\begin{array}{c}
       \varphi\\
       \psi
     \end{array}\right) \]
  and look for a solution of the form
  \[ \left(\begin{array}{c}
       \varphi\\
       \psi
     \end{array}\right) = \left(\begin{array}{c}
       Q_0 \widetilde{h}_1\\
       \rho \widetilde{h}_2
     \end{array}\right) . \]
  With $L_0 = \Delta - \frac{1}{r^2} + (1 - \rho^2)$ and $L_0 (\rho) = 0$, we
  have
  \[ \left(\begin{array}{c}
       Q_0 \left( \widetilde{h}_1'' + \left( 2 \frac{Q_0'}{Q_0} + \frac{1}{r}
       \right) \widetilde{h}_1' \right)\\
       \rho \left( \widetilde{h}_2'' + \left( 2 \frac{\rho'}{\rho} + \frac{1}{r}
       \right) \widetilde{h}_2' \right)
     \end{array}\right) = \bar{V}_{\lambda} \left(\begin{array}{c}
       Q_0 \widetilde{h}_1\\
       \rho \widetilde{h}_2
     \end{array}\right) . \]
  With the factorizations
\begin{align*}
& \widetilde{h}_1'' + \left( 2 \frac{Q_0'}{Q_0} + \frac{1}{r} \right)
     \widetilde{h}_1' = \frac{1}{Q_0^2 r} (Q_0^2 r \widetilde{h}_1')' \\
& \widetilde{h}_2'' + \left( 2 \frac{\rho'}{\rho} + \frac{1}{r} \right)
     \widetilde{h}_2' = \frac{1}{\rho^2 r} (\rho^2 r \widetilde{h}_2')',
\end{align*}
the first equation becomes
  \[ (Q_0^2 r \widetilde{h}_1')' = Q_0 r \bar{V}_{\lambda} \left(\begin{array}{c}
       Q_0 \widetilde{h}_1\\
       \rho \widetilde{h}_2
     \end{array}\right) \cdot \left(\begin{array}{c}
       1\\
       0
     \end{array}\right) \]
  and we choose (in order to have $\widetilde{h}_1 (0) = 1$)
  \[ \widetilde{h}_1 (r) = 1 + \int_0^r \frac{1}{Q_0^2 (t) t} \int_0^t s Q_0 (s)
     \left( \bar{V}_{\lambda} (s) \left(\begin{array}{c}
       Q_0 \widetilde{h}_1\\
       \rho \widetilde{h}_2
     \end{array}\right) (s) \right) \cdot \left(\begin{array}{c}
       1\\
       0
     \end{array}\right)\dd s\dd t. \]
  Similarly, we choose (in order to have $\widetilde{h}_2 (0) = 0$)
  \[ \widetilde{h}_2 (r) = \int_0^r \frac{1}{\rho^2 (t) t} \int_0^t s \rho (s)
     \left( \bar{V}_{\lambda} (s) \left(\begin{array}{c}
       Q_0 \widetilde{h}_1\\
       \rho \widetilde{h}_2
     \end{array}\right) (s) \right) . \left(\begin{array}{c}
       0\\
       1
     \end{array}\right) \dd s \dd t \]
  Now, we define $(h_1, h_2)$ such that
  \[ \left(\begin{array}{c}
       \varphi\\
       \psi
     \end{array}\right) = \left(\begin{array}{c}
       Q_0 \widetilde{h}_1\\
       \rho \widetilde{h}_2
     \end{array}\right) = Q_0 \left(\begin{array}{c}
       h_1\\
       h_2
     \end{array}\right), \]
  that is $h_1 = \widetilde{h}_1, h_2 = \frac{\rho}{Q_0} \widetilde{h}_2$, and we
  write the equation as the system on $h_1, h_2$:
  \[ \left(\begin{array}{c}
       h_1\\
       h_2
     \end{array}\right) = \left(\begin{array}{c}
       1\\
       0
     \end{array}\right) + \Theta \left(\begin{array}{c}
       h_1\\
       h_2
     \end{array}\right) \]
  where
  \begin{eqnarray*}
    \Theta \left(\begin{array}{c}
      h_1\\
      h_2
    \end{array}\right) & = & \left(\begin{array}{c}
      \Theta_1 (h_1, h_2)\\
      \Theta_2 (h_1, h_2)
    \end{array}\right)\\
    & = & \left(\begin{array}{c}
      \int_0^r \frac{1}{Q_0^2 (t) t} \int_0^t s Q_0^2 (s) \left(
      \bar{V}_{\lambda} (s) \left(\begin{array}{c}
        h_1\\
        h_2
      \end{array}\right) (s) \right) . \left(\begin{array}{c}
        1\\
        0
      \end{array}\right) \dd s \dd t\\
      \frac{\rho (r)}{Q_0 (r)} \int_0^r \frac{1}{\rho^2 (t) t} \int_0^t s \rho
      (s) Q_0 (s) \left( \bar{V}_{\lambda} (s) \left(\begin{array}{c}
        h_1\\
        h_2
      \end{array}\right) (s) \right) . \left(\begin{array}{c}
        0\\
        1
      \end{array}\right) \dd s \dd t
    \end{array}\right) .
  \end{eqnarray*}
  Since $Q_0 (r) \sim r$ near $r = 0$ and grows exponentially as $r \to + \infty$ while $\rho (r) \sim r$ when $r \to 0$ and
  $\rho (r) \rightarrow 1$ when $r \to + \infty$, we check that given
  $b_1, b_2 \in \mathbb{N}, b_3 \in \mathbb{Z}, c \in \{ 1, 2 \}$ we have
  \begin{equation}
    \int_0^t \frac{s^{b_1} \rho^{b_2} (s) Q_0^c (s)}{(1 + s)^{b_3}}\dd s
    \lesssim_{b_1, b_2, b_3, c} \frac{t^{b_1 + 1} \rho^{b_2} (s) Q_0^c (t)}{(1
    + t)^{b_3 + 1}} \label{swarm}
  \end{equation}
for any $t \geqslant 0$. By Lemma \ref{eqsec9}.3, we have
  \[ | \bar{V}_{\lambda} (s) | \lesssim \lambda^2 + \frac{\lambda}{(1 + s)^2},
  \]
  hence for $r \in [0, 8  y_0\xi^{- 1}]$, we have
  \begin{align*}
| \Theta_1 (h_1, h_2) (r) | & \lesssim \int_0^r \frac{1}{Q_0^2 (t) t} \int_0^t s Q_0^2 (s) \left(
    \lambda^2 + \frac{\lambda}{(1 + s)^2} \right) \dd s \dd t \| (h_1, h_2)
    \|_{L^{\infty} ([0, 8  y_0\xi^{- 1}])}\\
    & \lesssim \int_0^r \frac{1}{Q_0^2 (t) t} t Q_0^2 (t) \left( \lambda^2
    + \frac{\lambda}{(1 + t)^2} \right)\dd t \, \| (h_1, h_2) \|_{L^{\infty} ([0,
    8 y_0 \xi^{- 1}])}\\
    & \lesssim \int_0^r \left[ \lambda^2 + \frac{\lambda}{(1 + t)^2} \right] \dd t \, \| (h_1,
    h_2) \|_{L^{\infty} ([0, 8 y_0 \xi^{- 1}])}\\
    & \lesssim (r \lambda^2 + \lambda) \, \| (h_1, h_2) \|_{L^{\infty} ([0, 8 y_0
    \xi^{- 1}])}\\
    & \lesssim \lambda \| (h_1, h_2)\, \|_{L^{\infty} ([0, 8 y_0 \xi^{- 1}])}
  \end{align*}
  and we check similarly that
  \[ | (1 + r) \partial_r \Theta_1 (h_1, h_2) (r) | \lesssim \lambda\| (h_1, h_2)
     \|_{L^{\infty} ([0, 8 \xi^{- 1}])} \]
and
  \[ | (1 + r) \Theta_2 (h_1, h_2) (r) | + | (1 + r) \partial_r \Theta_2 (h_1,
     h_2) (r) | \lesssim \lambda \| (h_1, h_2) \|_{L^{\infty} ([0, 8 y_0 \xi^{-
     1}])} . \]
  With Lemma \ref{heimat}, we deduce that
  \[ \left\| \Theta \left(\begin{array}{c}
       h_1\\
       h_2
     \end{array}\right) \right\|_{C^1 ([0, 8 \xi^{- 1}])} \lesssim \lambda
     \left\| \left(\begin{array}{c}
       h_1\\
       h_2
     \end{array}\right) \right\|_{C^1 ([0, 8 y_0 \xi^{- 1}])} \]
  hence $\Theta$ is a contraction for the $C^1 ([0, 8 y_0 \xi^{- 1}])$ norm
  provided that $\lambda$ is small enough (depending on $y_0$).
  
  This completes the construction of $(h_1, h_2)$ on $[0, 8 y_0 \xi^{- 1}]$ with,
  for $k \in \{ 0, 1 \}$,
  \[ \left| \partial_r^k \left( \frac{1}{\lambda}\left( \left(\begin{array}{c}
       h_1\\
       h_2
     \end{array}\right) - \left(\begin{array}{c}
       1\\
       0
     \end{array}\right)\right)\right) \right| \lesssim_k (1 + r)^{- k} .
  \]
  Using the equation satisfied by $h_1$ and $h_2$, we check by induction that
  this still holds for any $k \in \mathbb{N}$.
  
  Now, we have
  \[  \frac{1}{\lambda}\left( \left(\begin{array}{c}
       h_1\\
       h_2
     \end{array}\right) - \left(\begin{array}{c}
       1\\
       0
     \end{array}\right) \right) = \frac{1}{\lambda } \Theta
     \left(\begin{array}{c}
       h_1\\
       h_2
     \end{array}\right) \]
  and the operator $\frac{1}{\lambda} \Theta $ is the same as $\Theta$ simply
  replacing $\bar{V}_{\lambda}$ by $\frac{\bar{V}_{\lambda}}{\lambda}$, which
  satisfies by Lemma \ref{eqsec9}.3 that for any $j \geqslant 1$,
  \[ \left| \partial^j_{\lambda} \left( \frac{\bar{V}_{\lambda}}{\lambda}
     \right) \right| \lesssim_j 1. \]
  We estimate then for $j \geqslant 1$ that
\begin{align*}
\left| \partial_{\lambda}^j \left( \frac{\Theta_1}{\lambda} \right) (h_1, h_2) (r) \right| & =  \int_0^r \frac{1}{Q_0^2 (t) t} \int_0^t s Q_0^2 (s) \left(
    \partial_{\lambda}^j \left( \frac{\bar{V}_{\lambda}}{\lambda} \right) (s)
    \left(\begin{array}{c}
      h_1\\
      h_2
    \end{array}\right) (s) \right) . \left(\begin{array}{c}
      1\\
      0
    \end{array}\right) \dd s \dd t\\
    & \lesssim  \int_0^r \frac{1}{Q_0^2 (t) t} \int_0^t s Q_0^2 (s) \dd s \dd t
    \| (h_1, h_2) \|_{L^{\infty} ([0, 8 y_0 \xi^{- 1}])}\\
    & \lesssim  \int_0^r \frac{1}{Q_0^2 (t) t} t Q_0^2 (t)\dd t \| (h_1, h_2)
    \|_{L^{\infty} ([0, 8  y_0\xi^{- 1}])}\\
    & \lesssim  (1 + r) \| (h_1, h_2) \|_{L^{\infty} ([0, 8 y_0 \xi^{- 1}])}
  \end{align*}
  and with similar estimates for $\partial_{\lambda}^j \left(
  \frac{\Theta_2}{\lambda} \right)$ and derivatives with respect to $r$, we can show the estimate on $\partial_r^k \partial_{\lambda}^j \left(
  \frac{1}{\lambda} \left[ \left(\begin{array}{c}
    h_1\\
    h_2
  \end{array}\right) - \left(\begin{array}{c}
    1\\
    0
  \end{array}\right) \right]\right) $.
  
  \
  
  We are left with the estimate of the function at $r = \nu y_0 \xi^{- 1}, \nu
  \in \{ 1/2, 4 \}$, and we need a refinement there.
  Similarly to the proof of Lemma \ref{heimat}, we check that the equation
  \[ L_0 - 2 \rho^2 - (\kappa^2 - 2) = 0 \]
  has a solution smooth near $r=0$ and belonging in $I_i (\kappa r) (1
  +\mathcal{W}_2 (r_\ast^+))$ since for $r \geqslant 1$,
  \[ L_0 - 2 \rho^2 - (\kappa^2 - 2) = \Delta - \kappa^2 +
     \frac{\kappa^2}{r^2} + O \left( \frac{1}{r^4} \right) . \]
  We denote by $\mathcal{I}_1$ this solution. Remark that since $\kappa^2-2>0$, we have $\mathcal{I}_1>0$ on $(0,\infty)$ by repeating exactly the proof of the positivity of $Q_0$ in Lemma \ref{lethargy}.
  
  Remark that by the definition of
  $\bar{V}_{\lambda}$, $\Phi_1, \Psi_1$ solves
  \[ \left(\begin{array}{c}
       (L_0 - 2 \rho^2 - (\kappa^2 - 2)) (\Phi_1)\\
       L_0 (\Psi_1)
     \end{array}\right) = \bar{\mathcal{V}}_{\lambda} \left(\begin{array}{c}
       \Phi_1\\
       \Psi_1
     \end{array}\right) . \]
  where
  \[ \bar{\mathcal{V}}_{\lambda} = \left(\begin{array}{cc}
       0 & - \lambda \frac{\rho^2 (r) - 1}{\langle \lambda \rangle}\\
       \frac{- \lambda}{\langle \lambda \rangle} (\rho^2 (r) - 1) & - 2 (1 -
       \langle \lambda \rangle) \frac{\rho^2 (r) - 1}{\langle \lambda \rangle}
       - \xi^2
     \end{array}\right) . \]
  This implies that if we write the solution as
  \[ \left(\begin{array}{c}
       \Phi_1\\
       \Psi_1
     \end{array}\right) =\mathcal{I}_1 \left(\begin{array}{c}
       \widetilde{h}_1\\
       \widetilde{h}_2
     \end{array}\right), \]
  then we have the implicit equation
  \[ \left(\begin{array}{c}
       \widetilde{h}_1\\
       \widetilde{h}_2
     \end{array}\right) = \left(\begin{array}{c}
       1\\
       0
     \end{array}\right) + \widetilde{\Theta} \left(\begin{array}{c}
       \widetilde{h}_1\\
       \widetilde{h}_2
     \end{array}\right) \]
  where
  \begin{eqnarray*}
    \widetilde{\Theta} \left(\begin{array}{c}
      \widetilde{h}_1\\
      \widetilde{h}_2
    \end{array}\right) & = & \left(\begin{array}{c}
      \int_0^r \frac{1}{\mathcal{I}_1^2 (t) t} \int_0^t s\mathcal{I}_1^2 (s)
      \left( \bar{\mathcal{V}}_{\lambda} (s) \left(\begin{array}{c}
        \widetilde{h}_1\\
        \widetilde{h}_2
      \end{array}\right) (s) \right) \cdot \left(\begin{array}{c}
        1\\
        0
      \end{array}\right) \dd s \dd t\\
      \frac{\rho (r)}{\mathcal{I}_1 (r)} \int_0^r \frac{1}{\rho^2 (t) t}
      \int_0^t s \rho (s) \mathcal{I}_1 (s) \left( \bar{\mathcal{V}}_{\lambda}
      (s) \left(\begin{array}{c}
        \widetilde{h}_1\\
        \widetilde{h}_2
      \end{array}\right) (s) \right) \cdot \left(\begin{array}{c}
        0\\
        1
      \end{array}\right) \dd s \dd t
    \end{array}\right) .
  \end{eqnarray*}
  We have
  \[ \left(\begin{array}{c}
       \widetilde{h}_1\\
       \widetilde{h}_2
     \end{array}\right) = \left(\begin{array}{c}
       1\\
       0
     \end{array}\right) + \widetilde{\Theta} \left(\begin{array}{c}
       1\\
       0
     \end{array}\right) + \widetilde{\Theta}^2 \left(\begin{array}{c}
       \widetilde{h}_1\\
       \widetilde{h}_2
     \end{array}\right) . \]
  The estimates coming from $\Theta^2 \left(\begin{array}{c}
    \widetilde{h}_1\\
    \widetilde{h}_2
  \end{array}\right)$ are easily checked. Now, we compute that
  \[ \widetilde{\Theta} \left(\begin{array}{c}
       1\\
       0
     \end{array}\right) = \left(\begin{array}{c}
       \int_0^r \frac{1}{\mathcal{I}_1^2 (t) t} \int_0^t s\mathcal{I}_1^2 (s)
       \left( \bar{\mathcal{V}}_{\lambda} (s) \left(\begin{array}{c}
         1\\
         0
       \end{array}\right) \right) \cdot \left(\begin{array}{c}
         1\\
         0
       \end{array}\right) \dd s \dd t\\
       \frac{\rho (r)}{\mathcal{I}_1 (r)} \int_0^r \frac{1}{\rho^2 (t) t}
       \int_0^t s \rho (s) \mathcal{I}_1 (s) \left(
       \bar{\mathcal{V}}_{\lambda} (s) \left(\begin{array}{c}
         1\\
         0
       \end{array}\right) \right) \cdot \left(\begin{array}{c}
         0\\
         1
       \end{array}\right) \dd s \dd t
     \end{array}\right), \]
  but now crucially,
  \[ \left( \bar{\mathcal{V}}_{\lambda} (s) \left(\begin{array}{c}
       1\\
       0
     \end{array}\right) \right) \cdot \left(\begin{array}{c}
       1\\
       0
     \end{array}\right) = 0. \]
  Furthermore,
  \[ \left( \bar{V}_{\lambda} (s) \left(\begin{array}{c}
       1\\
       0
     \end{array}\right) \right) \cdot \left(\begin{array}{c}
       0\\
       1
     \end{array}\right) = - 2 (1 - \langle \lambda \rangle) \frac{\rho^2 (r) -
     1}{\langle \lambda \rangle} - \xi^2, \]
  and we check that this implies that for all $j \in \mathbb{N}, \nu \in \{ \frac12,
  4 \}$, we have
  \[ \left| \partial_{\lambda}^j \left( \widetilde{\Theta} \left(\begin{array}{c}
       1\\
       0
     \end{array}\right) (\nu y_0 \xi^{- 1}) \right) \right| \lesssim_{j,y_0}
     \lambda^{\max (0, 2 - j)} . \]
  We can now conclude with \ Lemmas \ref{kast} and \ref{heimat}, following the
  end of the proof of Lemma \ref{alegant}.
\end{proof}

\end{document}
